\begin{notechapter}
\chapter{Measures}
\label{ch:measures}
\chaptersubtitle{%
  $\sigma$-algebras, Lebesgue measure and push-forwards}

\begin{chapterplan}
\begin{itemize}[nosep]
  \item Why the Riemann integral is not enough: limits of functions and
        expectations.
  \item $\sigma$-algebras and measures; continuity from below and from above,
        with the counterexample showing where finiteness matters; Borel sets;
        the $\pi$--$\lambda$ theorem.
  \item Lebesgue measure on $\R^n$: existence (\cref{ch:caratheodory}),
        uniqueness, invariance under translations (it is a measure on the affine
        space) and scaling by $\abs{\det A}$ under linear maps.
  \item Measurable maps and the push-forward of a measure.
  \item The vocabulary of probability: random variables and their laws.
\end{itemize}
\smallskip\noindent\textbf{Machine learning threads.}
\begin{itemize}[nosep]
  \item The law of a model's output is a push-forward.
  \item Generative models as push-forwards of a simple measure.
\end{itemize}
\end{chapterplan}
\end{notechapter}
